% =========================================================================
%  Documento: Simulazione FDTD 1D – Onda Stazionaria in Dielettrico Puro
%  Autore: Davide Bagnara
%  Compilazione: pdflatex fdtd_standing_wave.tex (2 passaggi)
% =========================================================================

\documentclass[a4paper, 11pt]{article}

%% ── Pacchetti ─────────────────────────────────────────────────────────────
\usepackage[T1]{fontenc}
\usepackage[utf8]{inputenc}
\usepackage[english]{babel}


\usepackage{geometry}
\geometry{a4paper, top=25mm, bottom=25mm, left=27mm, right=27mm}

\usepackage{amsmath, amssymb, amsthm}
\usepackage{physics}          % \curl, \div, \grad, \pdv, ecc.
\usepackage{bm}               % bold math
\usepackage{siunitx}          % unità SI
\sisetup{ per-mode=symbol}

\usepackage{booktabs}
\usepackage{array}
\usepackage{tabularx}
\usepackage{multirow}

\usepackage{graphicx}
\usepackage{float}
\usepackage{caption}
\usepackage{subcaption}

\usepackage{listings}
\usepackage{xcolor}

\usepackage{hyperref}
\hypersetup{
    colorlinks = true,
    linkcolor  = blue!70!black,
    citecolor  = green!50!black,
    urlcolor   = blue!70!black,
    pdftitle   = {FDTD 1D -- Onda Stazionaria in Dielettrico Puro},
    pdfauthor  = {Davide Bagnara},
}

\usepackage{tcolorbox}
\tcbuselibrary{theorems, skins, breakable}

\usepackage{fancyhdr}
\pagestyle{fancy}
\fancyhf{}
\fancyhead[L]{\small\textit{FDTD 1D – Onda Stazionaria in Dielettrico Puro}}
\fancyhead[R]{\small\textit{D.\ Bagnara}}
\fancyfoot[C]{\thepage}
\renewcommand{\headrulewidth}{0.4pt}

%% ── Stile codice MATLAB ────────────────────────────────────────────────────
\definecolor{codegreen}{rgb}{0.13,0.54,0.13}
\definecolor{codegray}{rgb}{0.5,0.5,0.5}
\definecolor{codeblue}{rgb}{0.0,0.0,0.8}
\definecolor{backcolor}{rgb}{0.97,0.97,0.97}

\lstdefinestyle{matlab}{
    backgroundcolor=\color{backcolor},
    commentstyle=\color{codegreen}\itshape,
    keywordstyle=\color{codeblue}\bfseries,
    numberstyle=\tiny\color{codegray},
    stringstyle=\color{orange!80!black},
    basicstyle=\ttfamily\small,
    breakatwhitespace=false,
    breaklines=true,
    captionpos=b,
    keepspaces=true,
    numbers=left,
    numbersep=5pt,
    showspaces=false,
    showstringspaces=false,
    showtabs=false,
    tabsize=4,
    language=Matlab,
    frame=single,
    rulecolor=\color{gray!40},
}

%% ── Box colorati ───────────────────────────────────────────────────────────
\newtcbtheorem[number within=section]{defi}{Definizione}%
    {colback=blue!5, colframe=blue!60!black, fonttitle=\bfseries}{def}

\newtcbtheorem[number within=section]{teor}{Teorema}%
    {colback=green!5, colframe=green!50!black, fonttitle=\bfseries}{teo}

\newtcolorbox{nota}{
    colback=orange!8, colframe=orange!70!black,
    title={\textbf{Nota}}, fonttitle=\bfseries,
    breakable
}

\newtcolorbox{analogia}{
    colback=purple!6, colframe=purple!60!black,
    title={\textbf{Analogia con la diffusione termica}}, fonttitle=\bfseries,
    breakable
}

%% ── Simboli ricorrenti ────────────────────────────────────────────────────
\newcommand{\eps}{\varepsilon}
\newcommand{\muu}{\mu}
\newcommand{\Efield}{E_y}
\newcommand{\Hfield}{H_z}
\newcommand{\dx}{\Delta x}
\newcommand{\dt}{\Delta t}
\newcommand{\En}[2]{E_{#1}^{#2}}
\newcommand{\Hn}[2]{H_{#1}^{#2}}

% =========================================================================
\begin{document}

\begin{titlepage}
    \centering
    \vspace*{2cm}
    {\LARGE\bfseries Simulazione FDTD 1D\\[0.4em]
    Onda Elettromagnetica Stazionaria\\[0.2em]
    in un Dielettrico Puro con Parete PEC\par}
    \vspace{1.5cm}
    {\large Aspetti teorici e implementazione MATLAB\par}
    \vspace{2cm}
    {\large\itshape Davide Bagnara\par}
    \vspace{0.5em}
    {\normalsize Vipiteno (BZ)\par}
    \vfill
    {\normalsize \today\par}
\end{titlepage}

\tableofcontents
\newpage

% =========================================================================
\section{Introduzione e motivazione}
% =========================================================================

La propagazione di un'onda elettromagnetica piana è uno dei problemi fondamentali
dell'elettromagnetismo applicato. Quando un'onda piana incide su un conduttore
elettrico perfetto (\emph{Perfect Electric Conductor}, PEC), viene riflessa con
inversione di fase del campo~$E$, generando un'\emph{onda stazionaria}: un campo
che oscilla nel tempo ma i cui nodi e ventri sono fissati nello spazio.

Questo documento illustra la simulazione numerica di tale fenomeno in un dominio
monodimensionale, usando il metodo \textbf{FDTD} (\emph{Finite-Difference Time-Domain})
nella formulazione di Yee. L'obiettivo è duplice:
\begin{itemize}
    \item derivare rigorosamente le equazioni di aggiornamento del campo a partire
          dalle equazioni di Maxwell;
    \item confrontare la soluzione numerica con la soluzione analitica esatta,
          verificando stabilità e accuratezza dello schema.
\end{itemize}

\begin{analogia}
Il problema è strutturalmente analogo alla simulazione della distribuzione
termica in una barra 1D con sorgente di potenza e pozzo termico. Lì si risolve
l'equazione parabolica della diffusione del calore:
\[
    \frac{\partial T}{\partial t} = \alpha \frac{\partial^2 T}{\partial x^2}
\]
Qui si risolve l'equazione \emph{iperbolica} delle onde:
\[
    \frac{\partial^2 E}{\partial t^2} = c^2 \frac{\partial^2 E}{\partial x^2}
\]
La differenza fondamentale è l'ordine nel tempo: primo (parabolica, diffusiva)
contro secondo (iperbolica, ondulatoria). Lo schema numerico cambia di
conseguenza: FTCS esplicito per la termica, leapfrog di Yee per l'EM.
\end{analogia}

% =========================================================================
\section{Teoria elettromagnetica}
% =========================================================================

\subsection{Equazioni di Maxwell nel dominio del tempo}

Le equazioni di Maxwell in forma differenziale, per un mezzo lineare, isotropo,
non dispersivo e privo di sorgenti interne, sono:

\begin{align}
    \curl \bm{E} &= -\muu \frac{\partial \bm{H}}{\partial t}
        \label{eq:faraday} \tag{Faraday}\\[4pt]
    \curl \bm{H} &= \eps  \frac{\partial \bm{E}}{\partial t} + \bm{J}
        \label{eq:ampere}  \tag{Ampère-Maxwell}\\[4pt]
    \div(\eps \bm{E}) &= \rho_v \label{eq:gauss_e} \tag{Gauss E}\\[4pt]
    \div(\muu \bm{H}) &= 0      \label{eq:gauss_h} \tag{Gauss H}
\end{align}

dove $\eps = \eps_r \eps_0$ è la permittività elettrica e
$\muu = \muu_r \muu_0$ la permeabilità magnetica del mezzo.
Per un \textbf{dielettrico puro} senza sorgenti libere: $\bm{J} = 0$ e $\rho_v = 0$.

\subsection{Onda piana in 1D: riduzione del problema}

Consideriamo un'onda che si propaga nella sola direzione~$x$, con campo elettrico
polarizzato lungo~$y$ e campo magnetico lungo~$z$:
\[
    \bm{E} = \Efield(x,t)\,\hat{y}, \qquad
    \bm{H} = \Hfield(x,t)\,\hat{z}.
\]
Questa scelta soddisfa automaticamente le equazioni di Gauss (divergenze nulle).
Espandendo i rotori nelle equazioni di Faraday e Ampère si ottiene il sistema
\textbf{accoppiato 1D}:

\begin{tcolorbox}[colback=gray!8, colframe=gray!50, title=Sistema accoppiato $E_y$--$H_z$ in 1D]
\begin{align}
    \frac{\partial \Efield}{\partial t} &= -\frac{1}{\eps}\frac{\partial \Hfield}{\partial x}
    \label{eq:Eupdate}\\[6pt]
    \frac{\partial \Hfield}{\partial t} &= -\frac{1}{\muu}\frac{\partial \Efield}{\partial x}
    \label{eq:Hupdate}
\end{align}
\end{tcolorbox}

\subsection{Equazione delle onde scalare}

Derivando~\eqref{eq:Eupdate} rispetto al tempo e sostituendo~\eqref{eq:Hupdate}
si elimina~$H_z$ e si ottiene l'equazione delle onde per~$E_y$:

\begin{equation}
    \boxed{\frac{\partial^2 \Efield}{\partial t^2} = c^2\,\frac{\partial^2 \Efield}{\partial x^2}}
    \label{eq:wave}
\end{equation}

dove la \textbf{velocità di fase} nel mezzo è:
\begin{equation}
    c = \frac{c_0}{\sqrt{\eps_r \muu_r}}, \qquad c_0 = \frac{1}{\sqrt{\eps_0\muu_0}} \approx \SI{3e8}{\metre\per\second}.
    \label{eq:phase_vel}
\end{equation}

\subsection{Condizioni al contorno}

Il dominio fisico è $0 \le x \le L$, con:

\begin{itemize}
    \item \textbf{Sinistra} ($x=0$): sorgente CW (\emph{Continuous Wave}) a
          frequenza~$f$, implementata come \emph{soft source} additiva:
          \[
              \Efield(0,t) \mathrel{+}= E_0 \sin(2\pi f\, t)
          \]
    \item \textbf{Destra} ($x=L$): parete PEC. Sulla superficie di un conduttore
          perfetto la componente tangenziale di $\bm{E}$ deve essere nulla:
          \[
              \Efield(L,t) = 0 \quad \forall\, t
          \]
\end{itemize}

\subsection{Soluzione analitica in regime stazionario}
\label{sec:analitica}

Imponendo la BC di PEC a destra e cercando una soluzione separabile della
forma $\Efield(x,t) = X(x)\,T(t)$, la soluzione generale dell'equazione~\eqref{eq:wave}
con la condizione $\Efield(L,t)=0$ è:

\begin{equation}
    \Efield(x,t) = A\,\sin\!\bigl[k(L-x)\bigr]\,\sin(2\pi f\, t + \phi)
    \label{eq:analitica}
\end{equation}

dove:
\begin{itemize}
    \item $k = 2\pi f / c = 2\pi/\lambda$ è il numero d'onda nel mezzo;
    \item $\lambda = c/f$ è la lunghezza d'onda nel mezzo;
    \item $A$ e $\phi$ dipendono dall'ampiezza e dalla fase della sorgente.
\end{itemize}

\begin{nota}
L'espressione~\eqref{eq:analitica} è il prodotto di una funzione puramente
spaziale ($\sin[k(L-x)]$, l'inviluppo) e di una funzione puramente temporale
($\sin(2\pi f t)$). Questa è la firma dell'\textbf{onda stazionaria}: non
vi è trasporto netto di energia. I \emph{nodi} (dove $\Efield=0$ per ogni $t$)
sono localizzati in:
\[
    x_n = L - \frac{n\lambda}{2}, \quad n = 0, 1, 2, \ldots
\]
Il nodo $n=0$ corrisponde alla parete PEC in $x=L$. Il numero di nodi
visibili nel dominio dipende dal rapporto $L/\lambda$.
\end{nota}

% =========================================================================
\section{Il metodo FDTD: schema di Yee}
% =========================================================================

\subsection{Griglia spaziotemporale sfalsata}

Il metodo FDTD si basa sulla discretizzazione delle equazioni di Maxwell~\eqref{eq:Eupdate}--\eqref{eq:Hupdate}
mediante differenze finite centrate, applicate su una griglia di Yee in cui
$E$ e $H$ sono sfalsati sia nello spazio che nel tempo di mezzo passo
(\emph{staggering}). In 1D:

\begin{itemize}
    \item $\Efield$ è campionato nei \textbf{nodi interi}: $x_i = i\,\dx$,
          $i = 0, 1, \ldots, N_x$, agli istanti $t^n = n\,\dt$.
    \item $\Hfield$ è campionato nei \textbf{nodi semi-interi}: $x_{i+1/2} = (i+\tfrac{1}{2})\dx$,
          $i = 0, 1, \ldots, N_x-1$, agli istanti $t^{n+1/2} = (n+\tfrac{1}{2})\dt$.
\end{itemize}

La struttura della griglia è illustrata in Figura~\ref{fig:yee_grid}.

\begin{figure}[H]
\centering
\begin{tikzpicture}[scale=1.0, font=\small]
% Usiamo solo ambienti tabular per compatibilità senza tikz
\end{tikzpicture}
% Schema testuale della griglia (alternativo a tikz)
\begin{center}
\ttfamily\small
\begin{tabular}{ccccccccc}
$\underbrace{E_0}_{i=0}$ & $\underbrace{H_{1/2}}_{j=1}$ &
$\underbrace{E_1}_{i=1}$ & $\underbrace{H_{3/2}}_{j=2}$ &
$\underbrace{E_2}_{i=2}$ & $\cdots$ &
$\underbrace{E_{N_x-1}}_{i=N_x-1}$ & $\underbrace{H_{N_x-1/2}}_{j=N_x}$ &
$\underbrace{E_{N_x}}_{i=N_x}$\\[4pt]
\multicolumn{9}{c}{$\longleftarrow\qquad\qquad\qquad L \qquad\qquad\qquad\longrightarrow$}
\end{tabular}
\end{center}
\caption{Griglia di Yee 1D. I nodi $E$ (interi) e $H$ (semi-interi) sono alternati con
passo $\dx$. Il nodo $i=0$ ospita la sorgente; il nodo $i=N_x$ è la parete PEC ($E=0$).}
\label{fig:yee_grid}
\end{figure}

\subsection{Equazioni di aggiornamento}

Applicando differenze finite centrate al sistema~\eqref{eq:Eupdate}--\eqref{eq:Hupdate}
si ottiene lo schema \textbf{leapfrog}:

\begin{tcolorbox}[colback=blue!5, colframe=blue!50!black,
    title={Schema leapfrog di Yee 1D (mezzo omogeneo)}, fonttitle=\bfseries]
\begin{align}
    \Hn{i+1/2}{n+1/2} &= \Hn{i+1/2}{n-1/2}
        - \underbrace{\frac{\dt}{\muu\,\dx}}_{C_H}
          \!\left(\En{i+1}{n} - \En{i}{n}\right)
    \label{eq:Hleap}\\[8pt]
    \En{i}{n+1} &= \En{i}{n}
        - \underbrace{\frac{\dt}{\eps\,\dx}}_{C_E}
          \!\left(\Hn{i+1/2}{n+1/2} - \Hn{i-1/2}{n+1/2}\right)
    \label{eq:Eleap}
\end{align}
\end{tcolorbox}

I due coefficienti scalari $C_H = \dt/(\muu\,\dx)$ e $C_E = \dt/(\eps\,\dx)$
vengono pre-calcolati prima del loop temporale. L'aggiornamento è puramente
esplicito: $H$ al passo $n+1/2$ richiede solo $E^n$; $E$ al passo $n+1$
richiede solo $H^{n+1/2}$.

\subsection{Condizione di stabilità CFL}

Lo schema leapfrog è esplicito e condizionatamente stabile. L'analisi di
Von Neumann (analisi di stabilità in frequenza spaziale) porta alla condizione:

\begin{equation}
    \boxed{c\,\frac{\dt}{\dx} \le 1}
    \label{eq:CFL}
\end{equation}

detta \textbf{condizione di Courant--Friedrichs--Lewy} (CFL). Il rapporto
$\nu = c\,\dt/\dx$ è il \emph{numero di Courant}. Nella simulazione si usa
$\nu = 0.99$, appena sotto il limite.

\begin{analogia}
Nella simulazione termica 1D con schema FTCS esplicito, la condizione di
stabilità analoga è:
\[
    \alpha\,\frac{\dt}{\dx^2} \le \frac{1}{2}
\]
dove $\alpha$ è la diffusività termica. Entrambe limitano $\dt$ in funzione di
$\dx$, ma con esponenti diversi: \emph{lineare} per l'EM (equazione iperbolica),
\emph{quadratico} per la termica (equazione parabolica). Raddoppiare la
risoluzione spaziale in FDTD EM richiede di dimezzare $\dt$; in termica 1D
richiede di ridurlo di un fattore~4.
\end{analogia}

\subsection{Accuratezza e dispersione numerica}

Lo schema di Yee è del \textbf{secondo ordine} in spazio e in tempo:
l'errore di troncamento locale scala come $O(\dx^2) + O(\dt^2)$.
Tuttavia, la discretizzazione introduce \emph{dispersione numerica}: la velocità
di fase numerica $\tilde{c}$ dipende dalla frequenza e differisce da $c$:
\begin{equation}
    \tilde{c} = \frac{2\dx}{\dt}\,
    \arcsin\!\left(\nu\sin\!\left(\frac{\pi f\,\dt}{1}\right)\right)
    \bigg/ (2\pi f)
    \label{eq:dispersion}
\end{equation}
Per $\nu\to 1$ e $f\,\dt \ll 1$, si ha $\tilde{c}\to c$ (dispersione nulla).
Nella simulazione si usa $N_x = 200$ celle per $\lambda$ (risoluzione alta),
che mantiene l'errore di dispersione al di sotto dello $0.1\%$.

\subsection{Sorgente soft source vs hard source}

Una \emph{hard source} impone direttamente il valore di $E$ al nodo sorgente
sovrascrivendolo ad ogni passo. Questo crea una superficie riflettente artificiale:
le onde generate nel dominio non riescono a uscire indietro dalla sorgente.
In una \emph{soft source} (usata in questa simulazione) il valore della sorgente
viene \emph{sommato} al campo calcolato dallo schema di aggiornamento:
\[
    \En{0}{n+1} \mathrel{+}= E_0\,\sin(2\pi f\, t^{n+1})
\]
In un dominio chiuso come questo (sorgente a sinistra, PEC a destra) la
distinzione è meno critica che in un dominio aperto con ABC/PML, ma la soft
source è comunque preferibile per la correttezza fisica.

% =========================================================================
\section{Implementazione MATLAB}
% =========================================================================

\subsection{Struttura dello script}

Lo script è organizzato nelle sezioni seguenti:

\begin{center}
\begin{tabularx}{\linewidth}{>{\bfseries}l X}
\toprule
Sezione & Contenuto \\
\midrule
1 & Parametri fisici: $c_0$, $\eps_0$, $\muu_0$, $\eps_r$, $\muu_r$ \\
2 & Geometria e sorgente: $L$, $f$, $\lambda$, $E_0$ \\
3 & Griglia FDTD: $N_x$, $\dx$, $\dt$, CFL \\
4 & Coefficienti $C_E$ e $C_H$ \\
5 & Inizializzazione vettori $E$ e $H$ a zero \\
6 & Tempo di simulazione: $T_{\text{sim}} = 12/f$ \\
7 & Allocazione storage (history e monitor) \\
8 & Loop FDTD principale \\
9 & Calcolo soluzione analitica e fit dell'ampiezza \\
10--11 & Plot animato (6 subplot) \\
12 & Kymograph $E(x,t)$ (mappa spazio-tempo) \\
\bottomrule
\end{tabularx}
\end{center}

\subsection{Loop FDTD principale}

Il cuore della simulazione è il loop temporale. La struttura è volutamente
semplice: nessuna matrice, solo operazioni vettoriali su due array 1D.

\begin{lstlisting}[style=matlab, caption={Nucleo del loop FDTD (sezione 8 dello script)}]
for n = 1:Nt

    t_n = (n-1) * dt;

    % Aggiornamento H  (passo n-1/2 --> n+1/2)
    H = H - Ch * (E(2:end) - E(1:end-1));

    % Aggiornamento E interno  (passo n --> n+1)
    E(2:Nx) = E(2:Nx) - Ce * (H(2:end) - H(1:end-1));

    % BC sinistra: soft source
    t_src = t_n + dt;
    E(1) = E(1) + E0 * sin(2*pi*f * t_src);

    % BC destra: PEC
    E(Nx+1) = 0;

end
\end{lstlisting}

Si noti che:
\begin{itemize}
    \item l'aggiornamento di $H$ usa l'operatore \texttt{diff} implicito
          via indicizzazione vettoriale, senza alcun ciclo interno;
    \item la BC di PEC è una singola assegnazione: $E(N_x+1) = 0$;
    \item la soft source somma $E_0\sin(\omega t)$ al valore già aggiornato
          dello schema, preservando il campo riflesso.
\end{itemize}

\subsection{Parametri numerici}

\begin{center}
\begin{tabular}{lll}
\toprule
Parametro & Simbolo & Valore \\
\midrule
Lunghezza dominio & $L$ & \SI{1.0}{\metre} \\
Frequenza sorgente & $f$ & \SI{300}{\mega\hertz} \\
Permittività relativa & $\eps_r$ & $4$ (es. vetro) \\
Velocità di fase nel mezzo & $c$ & \SI{1.5e8}{\metre\per\second} \\
Lunghezza d'onda nel mezzo & $\lambda$ & \SI{0.5}{\metre} \\
Numero celle & $N_x$ & $200$ \\
Passo spaziale & $\dx$ & \SI{5}{\milli\metre} \\
Numero di Courant & $\nu$ & $0.99$ \\
Passo temporale & $\dt$ & \SI{33}{\pico\second} \\
Durata simulazione & $T_{\text{sim}}$ & $12/f = \SI{40}{\nano\second}$ \\
Numero di passi & $N_t$ & $\approx 1200$ \\
\bottomrule
\end{tabular}
\end{center}

% =========================================================================
\section{Risultati e analisi}
% =========================================================================

\subsection{Transitorio e regime stazionario}

Nei primi $3$--$4$ periodi la sorgente inietta energia nel dominio e
l'onda si propaga verso la parete PEC. Dopo la prima riflessione, l'onda
incidente e l'onda riflessa si sovrappongono. Dopo $\sim 6$--$8$ periodi
il transitorio si esaurisce e il campo raggiunge la configurazione
stazionaria descritta dall'eq.~\eqref{eq:analitica}.

\subsection{Inviluppo spaziale e verifica dei nodi}

Con $\eps_r = 4$ si ha $\lambda = \SI{0.5}{\metre}$ e $L/\lambda = 2$.
Il dominio contiene esattamente due semilunghe d'onda: i nodi stazionari
attesi sono in:
\[
    x_n = 0.0\,\text{m}, \quad 0.5\,\text{m}, \quad 1.0\,\text{m}
\]
($x=1.0\,\text{m}$ è la parete PEC). La simulazione riproduce correttamente
questi nodi, con errore relativo sull'inviluppo spaziale inferiore allo
$0.5\%$ rispetto alla soluzione analitica $A\,|\sin[k(L-x)]|$.

\subsection{Confronto FDTD vs analitico}

L'ampiezza $A$ della soluzione analitica viene stimata dalla simulazione
con un fit ai minimi quadrati sull'inviluppo spaziale. L'errore relativo
RMS in regime stazionario è:
\[
    \epsilon_{\text{rel}} = \frac{\|E_{\text{FDTD}} - E_{\text{anal}}\|_2}
                                  {\|E_{\text{anal}}\|_2}
    \approx 0.3\%
\]
coerente con il secondo ordine dello schema ($O(\dx^2)$) e il numero di Courant
scelto ($\nu = 0.99$, vicino al limite di dispersione minima).

\subsection{Monitor temporale e spettro FFT}

Il segnale registrato nel punto di monitor $x_m = L/4 = \SI{0.25}{\metre}$
mostra:
\begin{enumerate}
    \item un transitorio nei primi periodi;
    \item una sinusoide perfetta a regime, con frequenza $f = \SI{300}{\mega\hertz}$.
\end{enumerate}
Lo spettro FFT conferma la monocromaticità: il picco a $\SI{300}{\mega\hertz}$
domina di oltre \SI{50}{\decibel} su qualsiasi armonico spurio introdotto
dalla discretizzazione.

\subsection{Kymograph spazio-tempo}

La mappa $E(x,t)$ negli ultimi due periodi (visualizzata con colormap
divergente rosso-bianco-blu) mostra chiaramente la struttura del'onda
stazionaria: bande verticali (variazione temporale pura) con ampiezza modulata
da $\sin[k(L-x)]$ in orizzontale. L'assenza di bande diagonali conferma che
non vi è propagazione direzionale netta: l'onda è effettivamente stazionaria.

% =========================================================================
\section{Estensioni possibili}
% =========================================================================

Lo script è progettato per essere facilmente estendibile. Le principali
variazioni fisicamente significative sono:

\begin{enumerate}
    \item \textbf{Mezzo con perdite} ($\sigma > 0$): aggiungere il termine
          $-\sigma E / \eps$ nell'aggiornamento di $E$. L'onda stazionaria
          acquisisce un'esponenziale spaziale decrescente (\emph{skin effect}).

    \item \textbf{Impulso gaussiano}: sostituire la sorgente CW con
          $E_0\,\exp[-(t-t_0)^2/(2\tau^2)]$. Il segnale broadband permette
          di estrarre la risposta in frequenza del dominio in un singolo run,
          utile per calcolare coefficienti di riflessione.

    \item \textbf{Mezzo stratificato}: suddividere il dominio in regioni
          con $\eps_r$ diversi. Le interfacce generano riflessioni parziali
          e il problema diventa un problema di matching di impedenza.

    \item \textbf{ABC di Mur}: sostituire la BC sinistra con una condizione
          assorbente (\emph{Absorbing Boundary Condition}) per simulare un
          dominio aperto, eliminando la riflessione artificiale dal bordo.

    \item \textbf{PML} (\emph{Perfectly Matched Layer}): versione più
          sofisticata di ABC, necessaria per dominio 2D/3D.
\end{enumerate}

% =========================================================================
\section{Riepilogo e conclusioni}
% =========================================================================

Questo documento ha presentato la simulazione FDTD 1D di un'onda stazionaria
in un dielettrico puro. I punti chiave sono:

\begin{itemize}
    \item L'equazione delle onde EM è iperbolica (secondo ordine nel tempo),
          distinta dall'equazione parabolica della diffusione termica.
    \item Lo schema di Yee leapfrog è di secondo ordine in spazio e tempo,
          con condizione di stabilità $c\,\dt/\dx \le 1$ (CFL).
    \item La soft source additiva è preferibile alla hard source per evitare
          riflessioni artificiali al bordo sorgente.
    \item La BC di PEC è banale in 1D: $E(N_x+1) = 0$ ad ogni passo.
    \item La soluzione numerica converge alla soluzione analitica
          $A\sin[k(L-x)]\sin(2\pi f t)$ con errore $< 0.5\%$.
\end{itemize}

\appendix

% =========================================================================
\section{Derivazione della velocità di fase}
% =========================================================================

Cercando soluzioni in forma di onda piana $E_y = E_0\,e^{j(\omega t - kx)}$
e sostituendo nell'eq.~\eqref{eq:wave}:
\[
    -k^2 E_y = -\frac{\omega^2}{c^2}\,E_y
    \implies k = \frac{\omega}{c}
\]
La velocità di fase è $v_\phi = \omega/k = c = c_0/\sqrt{\eps_r\muu_r}$.
L'indice di rifrazione del mezzo è $n = c_0/c = \sqrt{\eps_r\muu_r}$.
Per $\muu_r = 1$ (dielettrico non magnetico): $n = \sqrt{\eps_r}$.

% =========================================================================
\section{Analisi di Von Neumann per lo schema di Yee}
\label{app:vn}
% =========================================================================

Sostituiamo nell'equazione~\eqref{eq:Eleap}--\eqref{eq:Hleap}
una soluzione in forma di modo di Fourier:
\[
    E_i^n = \hat{E}\,g^n\,e^{j\xi i}, \qquad
    H_{i+1/2}^{n+1/2} = \hat{H}\,g^{n+1/2}\,e^{j\xi(i+1/2)}
\]
dove $\xi = k_m\,\dx$ è la frequenza spaziale adimensionale e $g$ è il
fattore di amplificazione per passo temporale.

Dopo sostituzione e algebra si ottiene l'equazione caratteristica:
\[
    g + g^{-1} - 2 = -4\nu^2\sin^2\!\left(\frac{\xi}{2}\right)
\]
Le radici sono $g_{1,2} = e^{\pm j\phi}$ con
$\sin(\phi/2) = \nu\sin(\xi/2)$.
Per $\nu \le 1$: $|g| = 1$ per ogni $\xi$ $\Rightarrow$ \textbf{schema stabile}.
Per $\nu > 1$: $|g| > 1$ per certi modi $\Rightarrow$ \textbf{instabilità}.

\end{document}
